\section{Fase qualitativa: seleção e roteiro de entrevistas}
\label{sec:qualitativa}

% comandos locais desta seção: cabeçalhos em negrito no lugar de
% subseções, para não pesar sumário nem bookmarks (xcolor já carregado
% no preâmbulo do principal).
\newcommand{\parte}[1]{\par\bigskip\noindent{\large\bfseries #1}\par\nopagebreak\smallskip}
\newcommand{\bloco}[1]{\par\medskip\noindent\textbf{#1}\par\nopagebreak\smallskip}
\newcommand{\achado}[1]{%
  \par\smallskip
  \noindent\begingroup\small\itshape
  \textcolor{black!70}{Achado que a pergunta testa: #1}%
  \endgroup\par\smallskip}

\begin{center}
\fbox{\parbox{0.92\textwidth}{\centering\bfseries\color{red}
SEÇÃO DESATUALIZADA --- SELEÇÃO DO PAINEL DE 314, PENDENTE DE
REEXECUÇÃO\\[2pt]
\normalfont\normalsize Os blocos abaixo citam hospitais já excluídos do
painel definitivo (289 CNES, ETAPA F de 15/07/2026), entre eles Darcy
Vargas (2071371), Boldrini (2081482) e o Centro Hospitalar do Sistema
Penitenciário (3001466). Aguardando reexecução da seleção antes de
qualquer leitura ou uso operacional.}}
\end{center}
\bigskip

\noindent\textbf{Uma nota sobre o estatuto desta seção.} O que se segue é
uma \textbf{sugestão}, construída por critérios puramente estatísticos:
os hospitais foram escolhidos pela sua posição nas distribuições de
eficiência, permanência e produção do painel, e as perguntas foram
derivadas dos achados quantitativos da fase de estimação. Nada disso
substitui o julgamento de campo. A decisão sobre \textbf{quais hospitais
de fato importam e quais perguntas fazer} cabe à experiência e ao
conhecimento institucional da pesquisadora principal, Priscilla Reinisch
Perdicaris, cuja leitura da rede hospitalar paulista é o fator decisivo
desta fase. As listas e os roteiros abaixo são apenas um ponto de partida
informado pelos dados, a ser confirmado, reordenado, enxugado ou ampliado
conforme esse julgamento.

\parte{Apresentação e método}

Esta seção operacionaliza a fase qualitativa da pesquisa sobre o efeito
do modelo institucional de gestão no desempenho dos hospitais SUS de São
Paulo, sobre o mesmo painel balanceado de 314 hospitais de 2015 a 2025.
A fase de estimação, já concluída, deixou quatro achados que as
entrevistas devem explicar por dentro. O primeiro é que o regime
operacional mais intenso e a eficiência técnica maior das OSS são traços
robustos da rede, presentes em todos os métodos. O segundo é que o tempo
médio de permanência cerca de 10\% menor nas OSS é evidência sugestiva de
variação dentro do hospital, apoiada em apenas cinco conversores. O
terceiro é que a mortalidade menor concentra-se na trajetória dos
conversores e não sobrevive ao controle sintético, o que recomenda
cautela e investigação de mecanismo. O quarto é que o faturamento real por
saída não difere entre categorias: a diferença de gestão, se existe, não
aparece no faturamento. \textit{Ressalva de leitura (regime de pagamento): hospitais OSS são custeados por contrato de gestão com orçamento global, não por reembolso de AIH --- o valor registrado por saída é, para eles, produção sem consequência financeira direta; a diferença de faturamento pode refletir o próprio regime de pagamento, a variável cujo efeito se busca medir, e não uso de recursos ou qualidade assistencial.} Toda leitura de mortalidade nesta fase usa a
complexidade estrutural como referência de perfil, nunca a complexidade
ponderada pela mortalidade, para evitar circularidade.

A seleção de entrevistados separa um \textbf{núcleo de doze entrevistas
previstas} de uma \textbf{reserva de entrevistas apenas sugeridas}, a
serem realizadas conforme o tempo e os recursos permitirem. Nenhum bloco
é obrigatório; os cinco conversores da administração Direta para OSS são
apenas os \textbf{prioritários}, por ancorarem toda a leitura causal do
documento. O núcleo, descrito a seguir, reúne os conversores e os
desviantes positivos e negativos de eficiência persistente no suporte
comum de porte e complexidade. A reserva sugerida, logo depois, reúne os
pares casados de contraste com estrutura parecida e eficiência divergente
e os casos de maior discordância entre DEA e SFA. O grupo Privado, com
apenas três hospitais, e os dezoito hospitais de longa permanência ficam
fora de todos os blocos comparativos, pelo tamanho e pelo perfil
assistencial não comparável. As entrevistas são semiestruturadas: o
roteiro-base comum vale para todos os entrevistados dos blocos de gestão,
e os módulos específicos por perfil acrescentam perguntas próprias de
cada grupo. Cada pergunta traz, em itálico, o achado quantitativo que ela
testa; o entrevistador deve conhecer esses achados, mas não deve
revelá-los ao entrevistado antes da resposta.

\parte{O núcleo: hospitais previstos para entrevista}

\bloco{Bloco a: conversores da Direta para OSS (prioritários)}
\begin{center}\small
\begin{tabular}{llllrr}
\toprule
CNES & Hospital & Município & Categoria & DEA & SFA \\
\midrule
2081695 & Conjunto Hospitalar Sorocaba & Sorocaba & OSS (conversor) & 0,560 & 0,999 \\
2078287 & Centro De Referencia Da Saude Da Mulhe & Sao Paulo & OSS (conversor) & 0,825 & 0,773 \\
2082225 & Hospital Katia De Souza Rodrigues Taip & Sao Paulo & OSS (conversor) & 0,820 & 0,906 \\
2091755 & Hospital Geral De Vila Penteado Dr Jos & Sao Paulo & OSS (conversor) & 0,581 & 0,882 \\
2750511 & Hospital Estadual Dr Odilo Antunes De  & Presidente Prudente & OSS (conversor) & 0,695 & 0,877 \\
\bottomrule
\end{tabular}
\end{center}
Os cinco conversores são o material mais valioso da pesquisa, pois
permitem observar o mesmo hospital antes e depois da mudança de gestão.
Suas trajetórias nos quatro indicadores já constam da
Seção~\ref{sec:conversores} (Figuras~\ref{fig:conv2081695}
a~\ref{fig:conv2750511b}), e por isso não se repetem aqui.

\noindent\textbf{2081695, Conjunto Hospitalar Sorocaba} (Sorocaba, interior; Grande Porte, faixa Barcelona 5, complexidade estrutural 675,0). Conversor da Direta para OSS em 2019. Antes e depois (medianas): TMP de 8,29 para 7,13 dias (queda de 14,0\%); mortalidade geral de 0,09 para 0,08 (queda de 14,8\%); faturamento por saída de R\$ 2.734,63 para R\$ 2.401,47 (queda de 12,2\%); saídas hospitalares de 13.159,50 para 16.243,00 (alta de 23,4\%).\par\medskip
\noindent\textbf{2078287, Centro de Referência da Saúde da Mulher, Pérola Byington} (São Paulo, capital; Médio Porte, faixa Barcelona 4, complexidade estrutural 266,0). Conversor da Direta para OSS em 2023. Antes e depois (medianas): TMP de 4,12 para 3,40 dias (queda de 17,5\%); mortalidade geral de 0,08 para 0,05 (queda de 39,5\%); faturamento por saída de R\$ 1.989,40 para R\$ 1.423,76 (queda de 28,4\%); saídas hospitalares de 5.583,00 para 14.036,00 (alta de 151,4\%).\par\medskip
\noindent\textbf{2082225, Hospital Kátia de Souza Rodrigues, Taipas} (São Paulo, capital; Grande Porte, faixa Barcelona 4, complexidade estrutural 380,0). Conversor da Direta para OSS em 2025, com apenas um ano posterior e sem leitura de tendência. Antes e depois (medianas): TMP de 5,41 para 5,35 dias (queda de 1,1\%); mortalidade geral de 0,08 para 0,06 (queda de 18,1\%); faturamento por saída de R\$ 869,18 para R\$ 977,07 (alta de 12,4\%); saídas hospitalares de 9.804,00 para 10.839,00 (alta de 10,6\%).\par\medskip
\noindent\textbf{2091755, Hospital Geral de Vila Penteado} (São Paulo, capital; Grande Porte, faixa Barcelona 4, complexidade estrutural 350,0). Conversor da Direta para OSS em 2025, com apenas um ano posterior e sem leitura de tendência. Antes e depois (medianas): TMP de 5,03 para 5,50 dias (alta de 9,3\%); mortalidade geral de 0,07 para 0,08 (alta de 11,1\%); faturamento por saída de R\$ 940,08 para R\$ 1.246,13 (alta de 32,6\%); saídas hospitalares de 8.350,50 para 7.783,00 (queda de 6,8\%).\par\medskip
\noindent\textbf{2750511, Hospital Estadual de Presidente Prudente} (Presidente Prudente, interior; Médio Porte, faixa Barcelona 2, complexidade estrutural 94,0). Conversor da Direta para OSS em 2025, com apenas um ano posterior e sem leitura de tendência. Antes e depois (medianas): TMP de 3,72 para 3,91 dias (alta de 5,4\%); mortalidade geral estável em 0,01 (alta de 13,4\%); faturamento por saída de R\$ 1.132,81 para R\$ 1.082,27 (queda de 4,5\%); saídas hospitalares de 3.675,00 para 3.513,00 (queda de 4,4\%).\par\medskip

\bloco{Bloco b: desviantes positivos (quartil superior persistente)}
\begin{center}\small
\begin{tabular}{llllrr}
\toprule
CNES & Hospital & Município & Categoria & DEA & SFA \\
\midrule
7573162 & Hospital Regional De Jundiai & Jundiai & OSS & 0,904 & 0,999 \\
6164366 & Hospital Estadual Americo Brasiliense & Americo Brasiliense & Direta & 0,807 & 0,886 \\
5718368 & Hosp Mun M Boi Mirim & Sao Paulo & Público Municipal & 0,867 & 0,962 \\
2076896 & Hospital Sao Luiz Gonzaga & Sao Paulo & Filantrópico & 0,927 & 0,925 \\
\bottomrule
\end{tabular}
\end{center}
Um desviante positivo por categoria, todos no suporte comum de porte e
complexidade, escolhidos por manterem eficiência alta ao longo dos onze
anos. A pergunta que eles respondem é quais práticas concretas sustentam
esse giro alto de forma persistente.

\noindent\textbf{7573162, Hospital Regional de Jundiaí} (Jundiaí, interior; Médio Porte, faixa Barcelona 3, complexidade estrutural 246,0). Desviante positivo da categoria OSS (quartil superior persistente de eficiência técnica em DEA e SFA): medianas de 0,904 no DEA e 0,999 na SFA entre 2015 e 2025, no suporte comum de faixa 3 e Médio Porte.\par\medskip
\noindent\textbf{6164366, Hospital Estadual Américo Brasiliense} (Américo Brasiliense, interior; Médio Porte, faixa Barcelona 3, complexidade estrutural 187,0). Desviante positivo da categoria Direta, por critério relaxado (quartil superior de DEA e acima da mediana de SFA): medianas de 0,807 no DEA e 0,886 na SFA entre 2015 e 2025, no suporte comum de faixa 3 e Médio Porte.\par\medskip
\noindent\textbf{5718368, Hospital Municipal M'Boi Mirim} (São Paulo, capital; Grande Porte, faixa Barcelona 4, complexidade estrutural 469,0). Desviante positivo da categoria Público Municipal (quartil superior persistente de eficiência técnica em DEA e SFA): medianas de 0,867 no DEA e 0,962 na SFA entre 2015 e 2025, no suporte comum de faixa 4 e Grande Porte.\par\medskip
\noindent\textbf{2076896, Hospital São Luiz Gonzaga} (São Paulo, capital; Grande Porte, faixa Barcelona 3, complexidade estrutural 241,0). Desviante positivo da categoria Filantrópico (quartil superior persistente de eficiência técnica em DEA e SFA): medianas de 0,927 no DEA e 0,925 na SFA entre 2015 e 2025, no suporte comum de faixa 3 e Grande Porte.\par\medskip

\begin{figure}[H]
\centering
\includegraphics[width=.95\textwidth]{fig_sel_b_7573162.png}
\caption{Hospital Regional de Jundiaí (CNES 7573162), desviante positivo
da rede OSS.}
\label{fig:sel7573162}
\end{figure}

\comoler{cada figura de seleção traz os quatro indicadores do hospital
(linha preta) contra a mediana da sua categoria (linha tracejada
colorida), de 2015 a 2025. Onde a linha preta corre abaixo da tracejada
em TMP, mortalidade e faturamento, ou acima em produção, o hospital supera a
referência da própria rede. Em Jundiaí, a produção corre bem acima da
mediana OSS e o TMP bem abaixo, a assinatura de um desviante positivo.}

\begin{figure}[H]
\centering
\includegraphics[width=.95\textwidth]{fig_sel_b_6164366.png}
\caption{Hospital Estadual Américo Brasiliense (CNES 6164366), desviante
positivo da rede Direta.}
\label{fig:sel6164366}
\end{figure}

\begin{figure}[H]
\centering
\includegraphics[width=.95\textwidth]{fig_sel_b_5718368.png}
\caption{Hospital Municipal M'Boi Mirim (CNES 5718368), desviante
positivo da rede Público Municipal.}
\label{fig:sel5718368}
\end{figure}

\begin{figure}[H]
\centering
\includegraphics[width=.95\textwidth]{fig_sel_b_2076896.png}
\caption{Hospital São Luiz Gonzaga (CNES 2076896), desviante positivo da
rede Filantrópico.}
\label{fig:sel2076896}
\end{figure}

\bloco{Bloco c: desviantes negativos (quartil inferior persistente)}
\begin{center}\small
\begin{tabular}{llllrr}
\toprule
CNES & Hospital & Município & Categoria & DEA & SFA \\
\midrule
2028840 & Instituto De Infectologia Emilio Ribas & Sao Paulo & Direta & 0,447 & 0,410 \\
2698471 & Secao Hospital Municipal Dr Arthur Dom & Santos & Público Municipal & 0,377 & 0,699 \\
2078465 & Real Sociedade Portuguesa De Beneficen & Campinas & Filantrópico & 0,229 & 0,464 \\
\bottomrule
\end{tabular}
\end{center}
O espelho do bloco anterior: um hospital por categoria no quartil
inferior persistente de eficiência, para ouvir quais restrições
impedem o giro, na visão de quem opera, antes de confrontá-las com os
dados do painel.

\noindent\textbf{2028840, Instituto de Infectologia Emílio Ribas} (São Paulo, capital; Grande Porte, faixa Barcelona 4, complexidade estrutural 459,0). Desviante negativo da categoria Direta (quartil inferior persistente de eficiência técnica em DEA e SFA): medianas de 0,447 no DEA e 0,410 na SFA entre 2015 e 2025, no suporte comum de faixa 4 e Grande Porte.\par\medskip
\noindent\textbf{2698471, Hospital Municipal Dr. Arthur Domingues Pinto} (Santos, interior; Médio Porte, faixa Barcelona 3, complexidade estrutural 122,0). Desviante negativo da categoria Público Municipal (quartil inferior persistente de eficiência técnica em DEA e SFA): medianas de 0,377 no DEA e 0,699 na SFA entre 2015 e 2025, no suporte comum de faixa 3 e Médio Porte.\par\medskip
\noindent\textbf{2078465, Real Sociedade Portuguesa de Beneficência} (Campinas, interior; Médio Porte, faixa Barcelona 3, complexidade estrutural 128,0). Desviante negativo da categoria Filantrópico (quartil inferior persistente de eficiência técnica em DEA e SFA): medianas de 0,229 no DEA e 0,464 na SFA entre 2015 e 2025, no suporte comum de faixa 3 e Médio Porte.\par\medskip

\begin{figure}[H]
\centering
\includegraphics[width=.95\textwidth]{fig_sel_c_2028840.png}
\caption{Instituto de Infectologia Emílio Ribas (CNES 2028840),
desviante negativo da rede Direta.}
\label{fig:sel2028840}
\end{figure}

\comoler{mesma leitura das figuras do bloco anterior, agora ao
contrário: nestes hospitais a linha preta tende a correr acima da
mediana da categoria em TMP e faturamento, ou abaixo em produção, marcando a
posição no quartil inferior de eficiência.}

\begin{figure}[H]
\centering
\includegraphics[width=.95\textwidth]{fig_sel_c_2698471.png}
\caption{Hospital Municipal Dr. Arthur Domingues Pinto (CNES 2698471),
desviante negativo da rede Público Municipal.}
\label{fig:sel2698471}
\end{figure}

\begin{figure}[H]
\centering
\includegraphics[width=.95\textwidth]{fig_sel_c_2078465.png}
\caption{Real Sociedade Portuguesa de Beneficência (CNES 2078465),
desviante negativo da rede Filantrópico.}
\label{fig:sel2078465}
\end{figure}

\parte{A reserva sugerida: entrevistas adicionais}

Os dois blocos a seguir não integram o núcleo previsto; são
\textbf{sugestões} a acrescentar caso o tempo de campo permita. Os pares
casados aprofundam o contraste da comparação principal do documento; os
casos de divergência entre métodos ajudam a entender o que os modelos
de eficiência não capturam.

\bloco{Bloco d: pares casados de contraste (sugeridos)}
\begin{center}\small
\begin{tabular}{llllrr}
\toprule
CNES & Hospital & Município & Categoria & DEA & SFA \\
\midrule
2792141 & Hospital Regional De Cotia & Cotia & OSS & 0,891 & 0,999 \\
2079410 & Complexo Hospitalar Padre Bento De Gua & Guarulhos & Direta & 0,528 & 0,795 \\
2096412 & Santa Casa De Misericordia De Jacarei & Jacarei & Filantrópico & 0,596 & 0,880 \\
2071371 & Hospital Infantil Darcy Vargas Uga Iii & Sao Paulo & Direta & 0,541 & 0,719 \\
2079232 & Hospital Santa Barbara & Santa Barbara D'Oeste & Filantrópico & 0,523 & 0,805 \\
\bottomrule
\end{tabular}
\end{center}
Duas âncoras OSS de alta eficiência, casadas cada uma a hospitais de
estrutura muito parecida mas eficiência bem menor. A ideia é entrevistar
os pares em espelho: mesma faixa Barcelona, mesmo porte, contagem de
leitos próxima, e ainda assim eficiência divergente. O que a gestão faz
de diferente?

\noindent\textbf{2792141, Hospital Regional de Cotia} (Cotia, interior; Médio Porte, faixa Barcelona 4, complexidade estrutural 330,0). Âncora OSS de alta eficiência (DEA 0,891, SFA 0,999) para os pares casados.\par\medskip
\noindent\textbf{2079410, Complexo Hospitalar Padre Bento de Guarulhos} (Guarulhos, interior; Médio Porte, faixa Barcelona 4, complexidade estrutural 327,0). Par casado da âncora OSS de Cotia (CNES 2792141): mesmo porte, faixa Barcelona 4, 145 leitos contra 145, mas eficiência divergente (DEA 0,528 contra 0,891; SFA 0,795 contra 0,999). Estrutura parecida sob gestão Direta.\par\medskip
\noindent\textbf{2096412, Santa Casa de Misericórdia de Jacareí} (Jacareí, interior; Médio Porte, faixa Barcelona 4, complexidade estrutural 312,0). Par casado da âncora OSS de Cotia (CNES 2792141): mesmo porte, faixa Barcelona 4, 143 leitos contra 145, mas eficiência divergente (DEA 0,596 contra 0,891; SFA 0,880 contra 0,999). Estrutura parecida sob gestão Filantrópico.\par\medskip
\noindent\textbf{2071371, Hospital Infantil Darcy Vargas} (São Paulo, capital; Médio Porte, faixa Barcelona 3, complexidade estrutural 222,0). Par casado da âncora OSS de Jundiaí (CNES 7573162): mesmo porte, faixa Barcelona 3, 116 leitos contra 136, mas eficiência divergente (DEA 0,541 contra 0,904; SFA 0,719 contra 0,999). Estrutura parecida sob gestão Direta.\par\medskip
\noindent\textbf{2079232, Hospital Santa Bárbara} (Santa Bárbara d'Oeste, interior; Médio Porte, faixa Barcelona 3, complexidade estrutural 247,0). Par casado da âncora OSS de Jundiaí (CNES 7573162): mesmo porte, faixa Barcelona 3, 137 leitos contra 136, mas eficiência divergente (DEA 0,523 contra 0,904; SFA 0,805 contra 0,999). Estrutura parecida sob gestão Filantrópico.\par\medskip

\begin{figure}[H]
\centering
\includegraphics[width=.95\textwidth]{fig_sel_d_2792141.png}
\caption{Hospital Regional de Cotia (CNES 2792141), âncora OSS de alta
eficiência dos pares casados.}
\label{fig:sel2792141}
\end{figure}

\comoler{a leitura é a mesma das figuras anteriores. Nos pares casados,
o proveito está em comparar esta figura da âncora OSS com a do seu par
de estrutura parecida nas páginas seguintes: onde as trajetórias se
descolam, mora a diferença que a entrevista busca explicar.}

\begin{figure}[H]
\centering
\includegraphics[width=.95\textwidth]{fig_sel_d_2079410.png}
\caption{Complexo Hospitalar Padre Bento de Guarulhos (CNES 2079410),
par Direta da âncora de Cotia.}
\label{fig:sel2079410}
\end{figure}

\begin{figure}[H]
\centering
\includegraphics[width=.95\textwidth]{fig_sel_d_2096412.png}
\caption{Santa Casa de Misericórdia de Jacareí (CNES 2096412), par
Filantrópico da âncora de Cotia.}
\label{fig:sel2096412}
\end{figure}

\begin{figure}[H]
\centering
\includegraphics[width=.95\textwidth]{fig_sel_d_2071371.png}
\caption{Hospital Infantil Darcy Vargas (CNES 2071371), par Direta da
âncora de Jundiaí.}
\label{fig:sel2071371}
\end{figure}

\begin{figure}[H]
\centering
\includegraphics[width=.95\textwidth]{fig_sel_d_2079232.png}
\caption{Hospital Santa Bárbara (CNES 2079232), par Filantrópico da
âncora de Jundiaí.}
\label{fig:sel2079232}
\end{figure}

\bloco{Bloco e: divergência entre DEA e SFA (sugeridos)}
\begin{center}\small
\begin{tabular}{llllrr}
\toprule
CNES & Hospital & Município & Categoria & DEA & SFA \\
\midrule
3001466 & Centro Hospitalar Do Sistema Penitenci & Sao Paulo & OSS & 0,393 & 0,999 \\
2081482 & Boldrini Campinas & Campinas & Filantrópico & 0,817 & 0,690 \\
2089785 & Hospital Do Rim E Hipertensao & Sao Paulo & Filantrópico & 0,774 & 0,651 \\
2077590 & Ibcc & Sao Paulo & Filantrópico & 0,737 & 0,447 \\
2079720 & Hospital Guilherme Alvaro Santos & Santos & OSS & 0,491 & 0,999 \\
\bottomrule
\end{tabular}
\end{center}
Hospitais em que os dois métodos de eficiência discordam fortemente, sinal
de composição de casos atípica, insumo mal medido ou produção heterogênea
que os modelos não capturam. A entrevista busca esse fator omitido.

\noindent\textbf{3001466, Centro Hospitalar do Sistema Penitenciário} (São Paulo, capital; Grande Porte, faixa Barcelona 4, complexidade estrutural 305,0). Divergência entre métodos: posição 277 no DEA contra 34 na SFA entre 305 CNES (SFA favorável, DEA desfavorável). Hospital de longa permanência.\par\medskip
\noindent\textbf{2081482, Boldrini Campinas} (Campinas, interior; Médio Porte, faixa Barcelona 3, complexidade estrutural 167,0). Divergência entre métodos: posição 24 no DEA contra 252 na SFA entre 305 CNES (DEA favorável, SFA desfavorável).\par\medskip
\noindent\textbf{2089785, Hospital do Rim e Hipertensão} (São Paulo, capital; Grande Porte, faixa Barcelona 3, complexidade estrutural 227,0). Divergência entre métodos: posição 42 no DEA contra 265 na SFA entre 305 CNES (DEA favorável, SFA desfavorável).\par\medskip
\noindent\textbf{2077590, IBCC} (São Paulo, capital; Grande Porte, faixa Barcelona 3, complexidade estrutural 176,0). Divergência entre métodos: posição 70 no DEA contra 290 na SFA entre 305 CNES (DEA favorável, SFA desfavorável).\par\medskip
\noindent\textbf{2079720, Hospital Guilherme Álvaro} (Santos, interior; Grande Porte, faixa Barcelona 4, complexidade estrutural 486,0). Divergência entre métodos: posição 241 no DEA contra 29 na SFA entre 305 CNES (SFA favorável, DEA desfavorável).\par\medskip

\parte{Roteiro-base comum aos blocos de gestão}

O roteiro-base cobre seis domínios e dura de cinquenta a sessenta
minutos. Vale para os entrevistados dos blocos a, b e c do núcleo e, se
realizados, dos pares casados sugeridos. As perguntas são abertas; os
itens entre colchetes são estímulos de aprofundamento para o
entrevistador.

\bloco{Governança e desenho institucional}
\begin{enumerate}[itemsep=2pt]
  \item Como se organiza a direção do hospital: quem decide o quê entre a diretoria local, a mantenedora ou OSS e a Secretaria? [organograma real e formal; conselhos; tempo de mandato da direção]
  \par\nopagebreak\achado{a eficiência técnica maior das OSS é traço robusto da rede, com DEA e SFA concordando; a hipótese é que o desenho de governança seja o canal.}
  \item Que decisões precisam de autorização externa ao hospital e quanto tempo levam, do pedido à resposta? [contratação, compra acima de um teto, abertura de leito]
  \par\nopagebreak\achado{o regime operacional mais intenso das OSS sugere menor fricção decisória; a pergunta mede a fricção diretamente.}
  \item Com que frequência a direção examina indicadores assistenciais e o que acontece quando um indicador piora? [rotina de reunião; quem responde pelo indicador]
  \par\nopagebreak\achado{desviantes positivos mantêm eficiência alta por onze anos; persistência sugere rotina gerencial, não sorte.}
\end{enumerate}

\bloco{Autonomia de recursos humanos e compras}
\begin{enumerate}[itemsep=2pt]
  \item Como o hospital contrata e desliga pessoal assistencial, e quanto tempo passa entre a necessidade identificada e o profissional trabalhando? [concurso e CLT; terceirização; banco de horas]
  \par\nopagebreak\achado{a produção das OSS é cerca de 45\% maior a estrutura comparável; a elasticidade de RH é o mecanismo candidato.}
  \item Como funcionam as compras e a manutenção de equipamentos: prazos típicos, o que mais trava? [licitação e compra direta; estoque; equipamento parado]
  \par\nopagebreak\achado{o faturamento real por saída não difere entre categorias; se as OSS produzem mais com faturamento igual, a diferença deve estar no uso dos insumos, não no preço.}
  \item Qual é a rotatividade das equipes médica e de enfermagem, e o que a explica? [vínculos múltiplos; plantões; salário relativo]
  \par\nopagebreak\achado{a estabilidade de equipe é explicação concorrente à gestão para a eficiência persistente dos desviantes positivos.}
\end{enumerate}

\bloco{Contrato de gestão e metas}
\begin{enumerate}[itemsep=2pt]
  \item Existe contrato, convênio ou instrumento com metas quantitativas? Quais indicadores, quem apura, e o que acontece quando a meta não é cumprida? [glosa; aditivo; renegociação]
  \par\nopagebreak\achado{o contrato de gestão é a diferença institucional formal entre OSS e administração direta; a queda de 10,2\% no TMP e a alta de 45\% na produção são os efeitos que ele teria de explicar.}
  \item As metas chegam ao cotidiano clínico? Dê um exemplo de decisão de ponta (alta, agenda cirúrgica, escala) que mudou por causa de meta. [quem traduz a meta para a equipe]
  \par\nopagebreak\achado{o TMP menor nas OSS é evidência dentro do hospital, apoiada em cinco conversores; a pergunta busca o elo micro entre contrato e permanência.}
\end{enumerate}

\bloco{Relação com a Secretaria e financiamento}
\begin{enumerate}[itemsep=2pt]
  \item Como é a interação com a Secretaria, estadual ou municipal: cobrança, apoio técnico, regulação de fila? Com que frequência? [câmaras técnicas; visitas; sistemas]
  \par\nopagebreak\achado{categorias públicas e OSS respondem à mesma Secretaria com resultados diferentes; a pergunta separa o vínculo formal da relação praticada.}
  \item Como o dinheiro chega e o que acontece numa crise de caixa? [teto MAC; incentivos; atraso de repasse; quem socorre]
  \par\nopagebreak\achado{o faturamento real por saída não difere entre categorias: se o financiamento é semelhante, a diferença de desempenho não é explicada por mais dinheiro.}
\end{enumerate}

\bloco{Gestão de leitos e fluxo do paciente}
\begin{enumerate}[itemsep=2pt]
  \item Descreva a rotina de altas: horário, critérios, quem decide, existe reunião diária de leitos? [alta assistida; previsão de alta na admissão]
  \par\nopagebreak\achado{o TMP cerca de 10\% menor nas OSS torna a rotina de gestão de altas o mecanismo mais direto que a entrevista pode verificar.}
  \item Quem controla a fila de internação e a ocupação do leito (NIR, regulação própria, CROSS)? Como se decide quem entra? [leito de retaguarda; internação social]
  \par\nopagebreak\achado{o giro alto dos desviantes positivos exige controle fino de entrada e saída; a pergunta identifica o instrumento.}
  \item Que parte da cirurgia migrou para o regime ambulatorial ou hospital-dia, e o que impediu de migrar mais? [centro cirúrgico; anestesia; retorno]
  \par\nopagebreak\achado{produção maior com TMP menor é compatível com substituição de internação por ambulatório; a pergunta testa essa composição.}
\end{enumerate}

\bloco{Resposta à pandemia (2020 e 2021)}
\begin{enumerate}[itemsep=2pt]
  \item O que mudou em leitos, equipes e fluxo em 2020 e 2021, e quanto da expansão foi física de fato e quanto foi apenas cadastro? [leitos de campanha; UTI improvisada; registro no CNES]
  \par\nopagebreak\achado{as séries de 2020 e 2021 têm extremos de ocupação de UTI e denominadores frágeis; a distinção entre leito real e leito cadastrado é decisiva para a leitura desses anos.}
  \item O que a pandemia deixou de permanente na operação do hospital? [telessaúde; protocolos; equipes]
  \par\nopagebreak\achado{a estimação de robustez sem 2020 e 2021 preserva os achados principais; a pergunta explora se houve mudança estrutural que os dados posteriores já refletem.}
\end{enumerate}

\parte{Módulos específicos por perfil}

\bloco{Módulo A: conversores da Direta para OSS (bloco a)}

Reconstituir a linha do tempo da transição: quando a conversão começou a
ser discutida, quando o contrato foi assinado, quando a OSS assumiu de
fato. Atenção metodológica: no Pérola Byington e em Sorocaba a melhora
dos indicadores começou \emph{antes} da conversão formal, e é
indispensável perguntar o que já estava em curso no período anterior
(intervenção da Secretaria, troca de direção, obra, mudança de perfil),
pois a resposta separa o efeito da gestão OSS de uma tendência
pré-existente.

\begin{enumerate}[itemsep=2pt]
  \item Conte a linha do tempo da transição para OSS: primeira discussão, decisão, assinatura, assunção. O que aconteceu no hospital em cada etapa? [greves; judicialização; transição de direção]
  \par\nopagebreak\achado{as datas de conversão adotadas (2019, 2023 e três em 2025) vêm do cadastro; a linha do tempo real pode antecedê-las e reinterpretar o estudo de eventos.}
  \item O que mudou concretamente no primeiro ano de OSS: rotinas, chefias, sistemas, escalas? O que não mudou? [primeiros noventa dias; o que a OSS trouxe pronto]
  \par\nopagebreak\achado{no estudo de eventos, o TMP cai já nos primeiros anos posteriores; a pergunta identifica qual prática explica a queda rápida.}
  \item Quanto da equipe assistencial e das chefias permaneceu após a conversão? [estatutários cedidos; recontratação CLT]
  \par\nopagebreak\achado{a continuidade de equipe distingue efeito de gestão de efeito de composição: se a equipe é a mesma, a mudança é de método, não de pessoas.}
  \item Quais metas o contrato de gestão fixou no início e como elas mudaram? Alguma meta era vista como inatingível? [metas de produção e de qualidade]
  \par\nopagebreak\achado{a alta de 45\% na produção e a queda de 10,2\% no TMP dentro do hospital fazem das metas contratuais a explicação institucional candidata.}
  \item (Somente Pérola Byington e Sorocaba) Os indicadores já vinham melhorando antes da conversão formal. O que estava acontecendo no hospital nos dois ou três anos anteriores? [gestor interino; plano de saneamento; antecipação da OSS]
  \par\nopagebreak\achado{a melhora pré-conversão observada nesses dois hospitais ameaça a leitura causal do estudo de eventos; a resposta decide entre antecipação, tendência ou seleção.}
\end{enumerate}

\bloco{Módulo B: desviantes positivos (bloco b)}

O objetivo é extrair práticas concretas e transferíveis por trás do giro
alto e da eficiência persistente, com ceticismo ativo: pedir sempre o
exemplo, o documento, a rotina, sem aceitar declaração de princípios.

\begin{enumerate}[itemsep=2pt]
  \item O hospital está no quartil superior de eficiência da rede há onze anos. Na sua leitura, quais três práticas explicam isso? Para cada uma: desde quando, quem implantou, o que custou? [pedir exemplos datados]
  \par\nopagebreak\achado{a persistência no quartil superior em DEA e SFA simultaneamente afasta artefato de método; deve haver prática real por trás.}
  \item Como funciona, passo a passo, a gestão de altas num dia típico? [reunião de leitos; previsão de alta; alta aos sábados]
  \par\nopagebreak\achado{giro alto com TMP baixo é a assinatura operacional dos desviantes positivos.}
  \item Que procedimentos o hospital faz em regime ambulatorial que vizinhos comparáveis internam? [cirurgia ambulatorial; hospital-dia; protocolos]
  \par\nopagebreak\achado{a produção alta com leitos comparáveis sugere substituição ambulatorial; a pergunta lista os procedimentos.}
  \item Como o hospital se relaciona com a regulação de vagas: aceita tudo que a central manda? Recusa o quê, e com que argumento? [seleção de casos; perfil de demanda]
  \par\nopagebreak\achado{a eficiência alta pode ser gestão ou seleção de pacientes mais leves; a pergunta confronta a hipótese de seleção, a ler junto com a complexidade estrutural do hospital.}
\end{enumerate}

\bloco{Módulo C: desviantes negativos (bloco c)}

Simétrico ao módulo B, mas centrado nas restrições percebidas. O tom não
é de auditoria: o objetivo é entender o que impede o giro, na visão de
quem opera, e confrontar depois com as restrições objetivas do painel.

\begin{enumerate}[itemsep=2pt]
  \item Quais são hoje os três maiores obstáculos para o hospital girar mais os leitos? [financiamento; RH; perfil de demanda; física do prédio]
  \par\nopagebreak\achado{o hospital está no quartil inferior persistente de DEA e SFA; a pergunta colhe a explicação nativa antes de confrontá-la com os dados.}
  \item Descreva o paciente típico que fica internado mais tempo do que precisaria. Por que ele não sai? [internação social; retaguarda; exame que não sai]
  \par\nopagebreak\achado{TMP alto com ocupação alta indica leito ocupado por paciente que não precisaria dele; a pergunta identifica o gargalo.}
  \item O financiamento cobre o custo do que o hospital produz? Onde aperta primeiro? [tabela SUS; folha; incentivos]
  \par\nopagebreak\achado{o faturamento real por saída não difere entre categorias na rede; se a restrição percebida é financeira, ela precisa ser conciliada com esse achado.}
  \item Se o hospital pudesse mudar uma única regra (de RH, compra ou contrato), qual mudaria e o que aconteceria? [pergunta projetiva]
  \par\nopagebreak\achado{a resposta revela qual margem de autonomia o gestor acredita que falta, comparável em espelho com o que as OSS dizem ter.}
\end{enumerate}

\bloco{Módulo D: pares casados (bloco d, sugeridos)}

As perguntas dos domínios de recursos humanos e compras, contrato e
metas, e leitos e fluxo do roteiro-base devem ser feitas de forma
espelhada e com a mesma ordem nos dois hospitais do par, para permitir
comparação direta das respostas. Acrescentar:

\begin{enumerate}[itemsep=2pt]
  \item Num dia típico, quem decide a abertura de um leito extra, a contratação de um plantonista e a compra urgente de um insumo, e em quanto tempo? [responder para os três casos]
  \par\nopagebreak\achado{o par tem faixa Barcelona, porte e leitos parecidos com eficiência divergente; a velocidade decisória é a diferença de gestão mais plausível com estrutura igual.}
  \item Como é montada a escala médica do fim de semana e o que acontece com as altas de sexta a domingo? [alta no fim de semana; plantão de sexta]
  \par\nopagebreak\achado{diferenças de TMP entre pares estruturalmente parecidos costumam se concentrar no fim de semana; pergunta espelhada nos dois hospitais do par.}
\end{enumerate}

\bloco{Módulo E: divergência entre DEA e SFA (bloco e, sugeridos)}

A divergência entre os dois métodos sugere algo que os modelos não
capturam: composição de casos atípica, insumo mal medido, produção
heterogênea. A entrevista busca esse fator omitido.

\begin{enumerate}[itemsep=2pt]
  \item O que este hospital faz que hospitais parecidos não fazem: serviços, perfis de paciente, papéis na rede que não aparecem em leitos e saídas? [ensino; referência regional; porta aberta e referenciada]
  \par\nopagebreak\achado{DEA e SFA discordam fortemente sobre este hospital; a explicação provável é dimensão de produto ou insumo que nenhum dos dois modelos mede.}
  \item Os números oficiais de leitos, produção e permanência do hospital refletem a operação real? Onde o registro distorce? [leitos bloqueados; produção não faturada]
  \par\nopagebreak\achado{a divergência entre métodos também pode ser artefato de medida; a pergunta verifica a qualidade do dado antes de interpretar o resto da entrevista.}
\end{enumerate}

\parte{Logística e esforço}

O núcleo previsto soma \textbf{doze entrevistas plenas}, com roteiro-base
de cinquenta a sessenta minutos mais o módulo específico de vinte a
trinta minutos (blocos a, b e c). A reserva sugerida acrescenta até dez
entrevistas (cinco pares casados no bloco d e cinco casos de divergência
no bloco e), a realizar conforme o tempo de campo permitir. Havendo
necessidade de priorizar, a ordem recomendada começa pelos conversores do
bloco a, segue pelos desviantes dos blocos b e c, cujo contraste entre
quartil superior e inferior é o coração da fase, e só então avança para
os pares casados do bloco d, mantendo ao menos um par completo por âncora
e preferindo o par Direta, que espelha a comparação principal da
estimação, e por último para o bloco e. Recomenda-se gravar com
consentimento, transcrever integralmente e codificar as respostas contra
a lista de achados deste roteiro, para que a fase qualitativa feche o
ciclo: cada mecanismo declarado nas entrevistas deve ser reexaminável no
painel. Toda a lista, no entanto, permanece subordinada ao ponto de
partida desta seção: a escolha final dos hospitais e das perguntas é da
pesquisadora principal.
